\documentclass[12pt]{article}
\usepackage{amsmath,amssymb,amsthm}
\usepackage[margin=1in]{geometry}

\newtheorem{exercise}{Exercise}
\title{Measure and Integral: Chapter 2 Exercises}
\date{}

\begin{document}
\maketitle

\begin{exercise}
Let $f(x)=x\sin(1/x)$ for $0<x\le 1$ and $f(0)=0$. Show that $f$ is bounded and continuous on $[0,1]$, but that $V[f;0,1]=+\infty$.
\end{exercise}

\begin{proof}
For $0<x\le1$,
\[
 |f(x)|=x|\sin(1/x)|\le x\le1,
\]
while $f(0)=0$; hence $f$ is bounded.  It is continuous on $(0,1]$ as a product and composition of continuous functions.  Moreover,
\[
 |f(x)-f(0)|=|f(x)|\le x\longrightarrow0
 \qquad(x\downarrow0),
\]
so it is continuous at $0$ as well.

To prove that its total variation is infinite, set
\[
 x_k=\frac{2}{(2k+1)\pi}\qquad(k=0,1,2,\ldots).
\]
Then $x_k\downarrow0$ and
\[
 \sin(1/x_k)=\sin\!\left(\frac{(2k+1)\pi}{2}\right)=(-1)^k,
 \qquad f(x_k)=(-1)^k x_k.
\]
For each $N\ge1$, use the partition
\[
 0<x_N<x_{N-1}<\cdots<x_1<x_0<1
\]
(with $0$ and $1$ included as endpoints).  Its variation sum contains the successive terms
\begin{align*}
 \sum_{k=0}^{N-1}|f(x_{k+1})-f(x_k)|
 &=\sum_{k=0}^{N-1}(x_{k+1}+x_k)\\
 &\ge\sum_{k=1}^{N}x_k
 =\frac{2}{\pi}\sum_{k=1}^{N}\frac{1}{2k+1}.
\end{align*}
The last partial sums are unbounded, since the series over the odd reciprocals diverges.  Thus variation sums for $f$ are arbitrarily large, and therefore
\[
 V[f;0,1]=+\infty.
\]
\end{proof}

\end{document}
